% !PRINTED 9
\weberpage{9}

%<<CONT
\weq{5}{\dfrac{(x_1x_2)(x_3x_4)}{(x_1x_3)(x_2x_4)} = \dfrac{(x_1'x_2')(x_3'x_4')}{(x_1'x_3')(x_2'x_4')}.}

Wenn die Funktion $f(x)$ in dem elliptischen Differential

\[
du = \dfrac{dx}{\sqrt{f(x)}}
\]

gegeben ist, so sind damit auch die $x_1$, $x_2$, $x_3$, $x_4$ gegeben, und durch die lineare Transformation (3) kann man das Differential auf eine andere Form bringen, bei der drei der Größe $x_1'$, $x_2'$, $x_3'$, $x_4'$ beliebig gegebene Werte haben. % [sic] printed "drei der Größe"; grammatically expected "drei der Größen"
Man erhält die \gesperrt{Normalform}, wenn man

\[
x' = z, \quad x_1' = 0, \quad x_2' = 1, \quad x_3' = \dfrac{1}{\varkappa^2}, \quad x_4' = \infty
\]

annimmt. Die Größe $\varkappa$ heißt bei \gesperrt{Legendre} der \gesperrt{Modul des elliptischen Differentials} und $\sqrt{1-\varkappa^2} = \varkappa'$ das \gesperrt{Komplement des Moduls}. Ordnet man also die Punkte in folgender Weise einander zu:

\[
\begin{array}{ccccc}
z, & 0, & 1, & \dfrac{1}{\varkappa^2}, & \infty, \\[6pt]
x, & x_1, & x_2, & x_3, & x_4,
\end{array}
\]

so ergeben sich aus der Gleichheit der Doppelverhältnisse leicht die Relationen:

\weq{6}{\begin{gathered}
z = \dfrac{(xx_1)(x_2x_4)}{(xx_4)(x_2x_1)}, \\
1 - z = \dfrac{(xx_2)(x_1x_4)}{(xx_4)(x_1x_2)}, \\
1 - \varkappa^2 z = \dfrac{(xx_3)(x_1x_4)}{(xx_4)(x_1x_3)},
\end{gathered}}

\weq{7}{dz = \dfrac{(x_4x_1)(x_4x_2)}{(x_2x_1)}\,\dfrac{dx}{(xx_4)^2},}

\weq{8}{\varkappa^2 = \dfrac{(x_3x_4)(x_1x_2)}{(x_1x_3)(x_2x_4)}, \quad \varkappa'^2 = \dfrac{(x_2x_3)(x_1x_4)}{(x_2x_4)(x_1x_3)}}

und aus (3):

\weq{9}{\dfrac{\sqrt{(x_3x_1)(x_4x_2)}\,dx}{\sqrt{-(x_1x)(x_2x)(x_3x)(x_4x)}} = \dfrac{dz}{\sqrt{z(1-z)(1-\varkappa^2z)}}.}

Nach § 2 wird $\varkappa^2$ ein positiver echter Bruch, wenn die $x_1$, $x_2$, $x_3$, $x_4$, reell sind und der Größe nach zyklisch aufeinander folgen. Dies gibt also acht verschiedene solche Transformationen,
%CONT>>
